\newif\ifglobal

%%%%%%%%%%%%%%%%%%%%%%%
%%% Compile to view %%%
%%%%%%%%%%%%%%%%%%%%%%%

%%%%%%%%%%%%%%%%%%%%%%%%%%%%%%%%%%%%%%%%%%%%%%%%%%%
%%% Readme:                                     %%%
%%% https://www.overleaf.com/read/bwdjvfjksvjy/ %%%
%%%%%%%%%%%%%%%%%%%%%%%%%%%%%%%%%%%%%%%%%%%%%%%%%%%

% >>> M?: marks a module setting number or important notice

% >>> M4: Set {./relative-path} to external files for this module <<<
\def \modulefiles {./00-BUSARCH}
% To add an external file, use:
% (Replace '---' with actual filename)
% '\input{\modulefiles/tables/---.tex}' for tables
% '\includegraphics{\modulefiles/figures/---}' for figures
% '\subfile{\modulefiles/---}' for register files


% >>> M1: This module is in global project? <<<
% 'YES' - uncomment; 'NO' - comment out
\globaltrue

\ifglobal
    \documentclass[main\_\_CN.tex]{subfiles}
   
\else
    % Font "FandolSong-Regular" used by ctex does not have GSUB table.
    % It causes warnings that do not affect compilation and can be ignored.
    \PassOptionsToPackage{quiet}{fontspec}

    \documentclass[openany, 10pt]{book}
   
    % >>> M2: In ./00-shared/preamble-trm-module, also find
    %         settings for visibility of todo notes, watermark <<< 
    \usepackage{preamble-trm-module}


    % Enable Chinese typesetting
    \usepackage{ctex}

    % >>> M3: Temporary labels in Chinese
    %         you can add your own labels there <<<
    \myexternaldocument{./00-shared/temp-labels-cn}

\fi %\ifglobal


\setcounter{secnumdepth}{4}

% Define formatting of chapters in text
\titleformat{\chapter}
  [display]
  {\LARGE\bfseries}
  {\color{AllportsColor}第\, \thechapter\, 章}
  {1em}
  {\color{AllportsColor}}
\titlespacing*{\chapter}{0pt}{0pt}{20pt}

% Define headers
\usepackage{fancyhdr}
\usepackage{lipsum}
\renewcommand{\chaptermark}[1]{\markboth{第\, \thechapter\, 章 \  \ #1}{}}
\fancypagestyle{plain}{%
  \fancyhead{}%
  %\fancyhead[R]{\Navigation}
  \fancyhead[L]{%
    \nouppercase{%
      \ %
      \normalsize\leftmark
    }%
  }%
}
\pagestyle{plain}

% >>> M5: If this module needs any updates in
% './00-shared/preamble-trm-module.sty'
% (include additional package, etc.), contact your module PM
% or create the file './00-shared/changelog.tex' make the
% necessary updates yourself and document them in this file <<<

% Make text left-aligned and allow latex to break pages naturally, comment out for full justification
\raggedright
\raggedbottom


\begin{document}


% Sets language into which pre-defined bits of text are translated
\selectlanguage{chinese}

% Do not delete this block
\ifglobal
% Add the following front matter if
% this module is outside of global project
\else
    \InputIfFileExists{readme}{}{}
    \listoftodos
    {\let\clearpage\relax \listofdonetodos}
    \tableofcontents
    %\listoftables
    %\listoffigures
\fi


%%%%%%%%%%%%%%%%%%%%%%%%%
%%% TRM Module Begins %%%
%%%%%%%%%%%%%%%%%%%%%%%%%

\hypertarget{busarch}{}
\chapter{总线架构}
\label{mod:busarch}

\chipname{} SoC 采用多层总线矩阵作为通信架构的核心，通过高效互连的结构连接多个主设备和从设备。该架构支持并发访问，有效缓解访问冲突，使高性能和低功耗的 RISC-V 核心、DMA 控制器、调试接口及各类外设能够在竞争最小的情况下同时运行。通过对数据流与系统性能进行优化，即使多个高速外设同时工作，总线矩阵也能够确保系统高效运转。

\chipname{} SoC 的主系统由以下部分组成：

\begin{itemize}
    \item 九个主设备 (master)：
        \begin{itemize}
            \item 高性能 (HP) 32 位 RISC-V 处理器的 I-bus、D-bus 和 AHB 总线
            \item 低功耗 (LP) 32 位 RISC-V 处理器
            \item TRACE
            \item TCM MEM MONITOR
            \item PSRAM MEM MONITOR
            \item GDMA
            \item SDIO SLAVE
        \end{itemize}
    \item 六个从设备 (slave)：
        \begin{itemize}
            \item ROM
            \item HP SRAM
            \item LP SRAM
            \item Cache
            \item HP\&LP 外设
            \item CPU 外设
        \end{itemize}
\end{itemize}

\chipname{} 的总线架构如下图所示。

\begin{figure}[H]
    \centering
    \includegraphics[width=1.0\textwidth]{\modulefiles/figures/bus-matrix_cn.pdf}
    \begin{minipage}{0.9\textwidth}
    \small
    注：\\
    1. 图中交叉点处的圆圈表示对应列的总线可以访问对应行的从设备。\\
    2. GDMA 上存在一条不经过总线矩阵、直接访问外部存储器的数据通路。\\
    3. CPU 访问外部存储器时必须通过 cache；通常情况下，cache 对用户而言是透明的。
    \end{minipage}
    \caption{\chipname{} 总线架构}
    \label{fig:bus-arch}
\end{figure}

图 \ref{fig:bus-arch} \textit{\nameref{fig:bus-arch}} 中的总线说明如下：

\begin{itemize}
    \item \textbf{HP CPU I-bus 总线}：连接高性能 (HP) 32 位 RISC-V 处理器的指令总线和总线矩阵，用于处理器取指。其可访问对象为包含代码的存储器（如 ROM、HP SRAM 或通过 MSPI 访问的外部存储器）。
    \item \textbf{HP CPU D-bus 总线}：连接高性能 (HP) 32 位 RISC-V 处理器的数据总线和总线矩阵，用于字面量加载和调试访问。其可访问对象为存储代码或数据的存储器（如 ROM、HP SRAM 或通过 MSPI 访问的外部存储器）。
    \item \textbf{HP CPU AHB 总线}：连接高性能 (HP) 32 位 RISC-V 处理器的 AHB 总线和总线矩阵，用于访问外设或 LP SRAM 中的数据。可访问对象包括 LP SRAM、HP\&LP 外设和 CPU 外设。
    \item \textbf{LP CPU 总线}：连接低功耗 (LP) 32 位 RISC-V 处理器的 AHB 总线和总线矩阵，供 LP CPU 访问内存和外设。可访问对象包括 HP SRAM、LP SRAM 和 HP\&LP 外设。
    \item \textbf{TRACE 总线}：连接 TRACE 模块和总线矩阵，用于执行调试操作。可访问对象包括 HP SRAM 和 LP SRAM。
    \item \textbf{TCM MEM MONITOR 总线}：连接 TCM MEM MONITOR 和总线矩阵，用于将 TCM MEM MONITOR 监控数据写入内存。可访问对象包括 HP SRAM 和 LP SRAM。
    \item \textbf{PSRAM MEM MONITOR 总线}：连接 PSRAM MEM MONITOR 和总线矩阵，用于将 PSRAM MEM MONITOR 记录的监控信息写入内存。可访问对象包括 HP SRAM 和 LP SRAM。
    \item \textbf{GDMA 总线}：提供一条连接至总线矩阵的线路和一条连接外部存储器的线路。GDMA 通过该总线进行高速数据传输，包括外设到内存、内存到外设、外设到外设以及内存到内存的传输。可访问对象包括 HP SRAM、LP SRAM 和外部存储器。针对 GDMA 访问外部存储器的场景，为确保数据吞吐率，建议尽可能设置较长的 burst 长度。
    \item \textbf{SDIO SLAVE 总线}：连接 SDIO SLAVE 和总线矩阵，用于访问内存或外设中的数据。可访问对象包括 HP SRAM、LP SRAM、HP\&LP 外设和 CPU 外设。
\end{itemize}

\begin{tiplisting}
通过 TRACE、TCM MEM MONITOR 或 PSRAM MEM MONITOR 访问 LP SRAM 的速度较慢，且 LP SRAM 空间较小，因此不建议使用这些路径访问 LP SRAM。
\end{tiplisting}

\end{document}
